\documentclass[11pt,letterpaper]{article}
\usepackage[T1]{fontenc}
\usepackage{lmodern}
\pdfmapfile{+lm.map}
\pdfmapfile{+symbols.map}
\pdfmapfile{+cm.map}
\pdfmapfile{+cmextra.map}
\usepackage[margin=0.82in]{geometry}
\usepackage{amsmath,amssymb,amsthm,mathtools}
\usepackage{booktabs,longtable,array,multicol}
\usepackage{xcolor,microtype}
\usepackage{xurl}
\usepackage[colorlinks=true,linkcolor=blue!45!black,urlcolor=blue!50!black]{hyperref}
\usepackage{fancyhdr,listings}
\definecolor{ink}{HTML}{16354D}
\hypersetup{pdftitle={PA Consistency Machines: Comparative Audit and Reduction Experiments},pdfauthor={Computational research report}}
\pagestyle{fancy}\fancyhf{}
\fancyhead[L]{\small\textcolor{ink}{PA consistency machines}}
\fancyhead[R]{\small Comparative audit and reductions}
\fancyfoot[C]{\thepage}
\setlength{\headheight}{14pt}
\setlength{\parskip}{0.35em}
\setlength{\parindent}{0pt}
\setlength{\emergencystretch}{2em}
\lstset{basicstyle=\ttfamily\footnotesize,breaklines=true,columns=fullflexible,frame=single,rulecolor=\color{black!20},backgroundcolor=\color{black!2},showstringspaces=false}
\newcommand{\PA}{\mathsf{PA}}
\newcommand{\Con}{\operatorname{Con}}
\newcommand{\Halt}{\operatorname{Halt}}
\newcommand{\NH}{\operatorname{Nonhalt}}
\begin{document}
\begin{titlepage}
\vspace*{0.5in}
{\color{ink}\Huge\bfseries PA Consistency Machines\par}
\vspace{0.25in}
{\Large A comparative audit and verified reduction experiments\par}
\vspace{0.2in}
{\large The supplied 209-, 208-, and 196-state constructions\par}
\vspace{0.3in}
{\large 13 November 2026\par}
\vspace{0.4in}
\hrule\vspace{0.2in}
\textbf{Result status.} No verified machine with fewer than 196 working states
was obtained. This report does not claim that 196 is a lower bound. It gives
an exact comparison of the supplied designs, two verified 196-state
intermediates, a new logical reduction of the source program, and reproducible
certificates distinguishing successful simulations from rejected smaller
candidates.

\textbf{Verified machine results.} A new rebracketing retains all 3,812 source
instructions and passes all 4,959 symbolic contracts. A second variant
transfers the $h$-only acceptance convention into PA196, has 3,811 source
instructions, and passes all 4,958 contracts. Independent literal interpreters
verify complete blank-tape startup configurations for both.

\textbf{Mathematical contributions.} A syntactic interpretation supplies the
PA-internal converse needed for the full consistency equivalence. A new
integrated equality-witness gate removes a primitive logical head and 67
source instruction sites, reducing the program from 3,812 to 3,745 sites.
Its best verified implementation has 203 working states. Achieving a physical
state count below 196 remains an additional obligation.

\textbf{Proof status.} The machine certificates are independently replayable
Python calculations justified by the published unbounded tape-engine lemmas.
The new logical and PA-internal arguments are mathematical proofs. This work
does not claim a new completed or executed Lean formalization.
\vfill
\small The accompanying ZIP contains this PDF, its complete TeX sources, literal
transition tables, source-site data, proof and verification scripts, source
revision locks, licenses, and selected search receipts.\par
\end{titlepage}
\tableofcontents
\newpage

\section{The target and the result obtained}
The requested target is a deterministic, one-tape, two-symbol machine $M$
with fewer than 196 working states, starting on the all-zero tape, for which
\[
       \PA\vdash \NH(M)\ \longleftrightarrow\ \Con(\PA).
\]
The halting control label is additional to the working-state count. The
meaning of $\Con(\PA)$ is the ordinary arithmetized absence of a PA proof of
contradiction, with the proof presentation fixed as in the cited development.
Changing to another ordinary proof presentation requires the usual
PA-provable primitive-recursive proof translations.

The strict state-count target remains unmet in this work. The positive
results below have complete, finite artifacts; a table with a nominal count
of 195 is expressly rejected because it fails a required simulation and
admits a four-step loop on a legitimate source frame. No inference of
nonhalting is drawn from merely running a machine for a long time.

\begin{center}
\small
\begin{tabular}{@{}lrrrr@{}}
\toprule
Verified table & States & Source sites & Contracts & Startup steps\\
\midrule
Published PA196 &196&3,812&4,959&58,154\\
New rebracketing &196&3,812&4,959&58,154\\
New $h$-only variant &196&3,811&4,958&57,922\\
Integrated witness, initial &206&3,745&4,872&59,164\\
Integrated witness, compressed &203&3,745&4,872&58,154\\
\bottomrule
\end{tabular}
\end{center}
The source-site numbers count primitive register instructions. The contract
numbers additionally distinguish the zero and positive cases of each
decrement. The startup numbers count literal binary tape transitions.
These three units are different and are not interchanged in state estimates.

The comparison in Section~\ref{sec:comparison-audit} explains why the
one-state U208 deletion cannot simply be applied to PA196. The unbounded
verification method and the new rebracketed table are given next. The later
sections establish the logical results and document the rejected reductions.

\input{comparison/audit/comparison_chapter.tex}

\section{Verification at the literal-machine boundary}
\label{sec:verification}
\subsection{Source frames and the complete register file}
The verifier independently transcribes the symbolic operations of the
published \texttt{M6/Runtime.lean}. A clean frame has control state 75, a
one immediately to the left of the head, and right-hand tape
\[
 p\,0^{\,n+48-|p|}\,1\,\operatorname{File}(v_0,\ldots,v_{k-1}),
 \qquad n\in\mathbb N,
\]
where $p$ is a source address, $k\geq9$, and
\[
 \operatorname{File}(\vec v)=
  0\,1^{v_0+1}\,0\,1^{v_1+1}\cdots0\,1^{v_{k-1}+1}.
\]
All remaining tape cells are zero. The logical register numbers $0,\ldots,8$
are stored in physical positions
\[
                  [6,2,0,4,3,7,8,5,1].
\]
The symbolic padding is an affine expression in the arbitrary natural $n$.
The register file is one opaque term, with a recorded operation and marked
register where appropriate. It is not replaced by a small test vector.

Four unbounded summaries cover marking the selected register, increment,
zero decrement, and positive decrement. Their justification is the existing
native tape-engine proofs. The verifier checks all 45 literal transition
rows on which those summaries depend. In particular the increment summary
includes the appended spare register with value one; changing that spare
value without adapting the engine proof is not an allowed optimization in
these positive certificates.

\subsection{What each finite contract proves}
For every source address the checker starts from its symbolic clean frame
and follows literal transition rows, finite scan summaries, and the four
unbounded register summaries. It must reach the prescribed next clean frame,
with the exact complete-file effect and a positive number of literal steps.
A halt instruction must reach the designated halting label. Decrement has
separate contracts under $v_r=0$ and $v_r>0$.

There are 1,147 decrement sites in the published source, hence
$3,812+1,147=4,959$ contracts. The check includes arithmetic-loop bodies and
control transfers as well as the straight-line formula constructors.
Unfolding a symbolic derivation and applying the engine lemmas gives a
finite positive simulation for every admissible natural padding and every
complete register file. The mathematical conclusion is therefore universal;
it is stronger than a collection of bounded numerical trials.

The release verifier uses explicit exceptions at its acceptance gates and
works with Python assertions disabled. The underlying source parser is
checked against pinned input hashes and the explicit graph identities.
This remains a small, reviewable certificate-replay implementation rather
than a replacement for a proof assistant's trusted kernel.

\subsection{Startup is checked independently}
Symbolic contracts begin at clean source frames, whereas the task starts on
a blank tape. A separate interpreter represents the entire tape by the set
of its nonzero cells and executes one literal transition at a time. It
imports neither the symbolic runtime nor its register summaries.
At the claimed startup step it compares the \emph{whole} tape, state, and
relative head position with the clean source-25 frame. The initial physical
file is
\[
       [0,0,0,0,0,0,2,0,0]\mathbin{+\!+}[1]^{26}.
\]
For both the published and the rebracketed table this occurs at step 58,154,
with absolute head position $-2$ and state 75. For the $h$-only table it occurs
at step 57,922. The zero extension of an address is part of its padding, so
the differing printed source-25 path of the $h$-only table does not hide a
tape discrepancy.

Positive step lengths, deterministic execution, and the one halting label
transport the source program's halting behavior to the literal blank-tape
machine. These transport arguments can be formalized in PA by bounded
induction and primitive-recursive trace construction. Section~\ref{sec:internal-converse}
spells out the additional syntactic argument needed for the consistency
biconditional.

\section{A structurally different verified 196-state machine}
\label{sec:rebracketed}
\subsection{Construction}
The AX body is a sequence over the following macro alphabet. The single-letter
names are opaque compression tokens: in particular token $\mathsf A$ copies
the register called $D$ in the arithmetic program, and token $\mathsf D$
copies the register called $A$. The mapping, not an intuitive interpretation
of the letters, defines the construction.
\begin{center}
\begin{tabular}{@{}clcl@{}}
\toprule
Token&Original source node&Token&Original source node\\
\midrule
$\mathsf P$&\texttt{node49}&$\mathsf H$&\texttt{node58}\\
$\mathsf A$&\texttt{node55}&$\mathsf D$&\texttt{node70}\\
$\mathsf C$&\texttt{node65}&$\mathsf R$&\texttt{node72}\\
$\mathsf S$&\texttt{node75}&$\mathsf G$&\texttt{node142}\\
$\mathsf K$&\texttt{node150}&$\mathsf i$&\texttt{inc 1}\\
$\mathsf B$&\texttt{breakTo 1}&&\\
\bottomrule
\end{tabular}
\end{center}
The search also exposes the closed copy suffixes
$\mathsf h=\texttt{node57}$, $\mathsf a=\texttt{node54}$, and
$\mathsf d=\texttt{node69}$, with
$\mathsf H=(\mathsf P,\mathsf h)$,
$\mathsf A=(\mathsf P,\mathsf a)$, and
$\mathsf D=(\mathsf P,\mathsf d)$.
Reassociations are accepted only after reconstructing every primitive
instruction's ordinal destinations and checking equality with the original
3,812-instruction control graph.

The selected grammar is
\path{verification/open_physical_61_201240.json}.
The ``61'' in this historical search filename is an \emph{obsolete optimistic
search score}; it is not the certified state saving. Its first fold proposal
contained an incompatible pair. The accepted certificate uses the compatible
six pairs below and has 196 states. This potentially confusing filename is
retained only so that the research trail remains reproducible.

Compiling its distinct new readers gives 202 working states, comprising the
unchanged 134-state exterior and 68 new pure readers. Six compatible
predecessor-sensitive identifications reduce 202 to 196. Finally a bijection
maps the surviving new reader labels to $134,\ldots,195$, preserving all labels
$0,\ldots,133$ and the halt label 196.

\subsection{The six finite local identifications}
In the temporary pure-reader numbering, the entire fold certificate is
\begin{center}
\begin{tabular}{@{}rrrr@{}}
\toprule
Kept&Deleted&Helper for read 0&Helper for read 1\\
\midrule
1004&1013&1039&--\\
1006&1028&1035&1023\\
1020&1003&1018&--\\
1029&1009&--&1057\\
1044&1024&1043&--\\
1048&1049&1008&--\\
\bottomrule
\end{tabular}
\end{center}
Every original entry to the kept state has predecessor bit zero, and every
original entry to the deleted state has predecessor bit one. A dash means
that the two original rows already agree. On a differing read bit $b$, the
shared state preserves $b$, moves left, and enters the stated helper. The
helper reads the predecessor bit, returns right, and selects a donor that
performs the original transition.

All 24 local cases are checked \emph{after all six identifications have been
applied}. Each takes one or three literal transitions, touches only the two
specified cells, preserves their contents, and finishes at the identified
old endpoint one cell to the right. This simultaneous check matters: an
individually suitable helper may cease to work when another selected fold
changes it. Ordinary unchanged rows commute with the identification.

The relabeling verifier checks all 392 rows, including halt targets. All 196
working states are graph reachable. The resulting canonical graph is
structurally different from the published PA196 graph, so this is more than
a relabeling of the published table.

\subsection{The concrete artifact and its guarantees}
The complete dense table is
\path{comparison/audit/rebracketed196_dense.tm}; it is also printed in
Appendix~\ref{app:table}. Its SHA-256 is
\begin{quote}\small
\nolinkurl{e212e359837cff6d1b3ab665855b0e9b4a2b471d23e77db447680d494d943b80}.
\end{quote}
It has start state 0, working states $0,\ldots,195$, and halt label 196.
The longest source path decreases from 38 to 37. Its 4,959 contracts have a
maximum symbolic execution length of 515 and a total of 1,589,112 symbolic
steps, compared with 1,619,679 for the published table. These symbolic step
counts are verifier costs, not universal bounds on literal tape runtimes.
The startup remains exactly 58,154 literal steps.

Since the primitive source program is unchanged, the all-natural contracts
and startup prove halting equivalence with PA196. Together with the
mathematical internal converse below, they give
\[
 \PA\vdash \NH(M_{\mathrm{rebracketed}})\leftrightarrow\Con(\PA).
\]
The state count is still 196; the construction supplies a verified different
starting point for further reductions.

\input{comparison/audit/h_only_shift_note.tex}
\input{math_reduction/internal_converse.tex}
\input{math_reduction/integrated_gate_reduction.tex}

\subsection{Verified implementations of the integrated gate}
The new logical source has a complete literal implementation in
\path{comparison/audit/integrated_gate206_dense.tm}.
It uses working states $0,\ldots,205$, start state 0, and halt label 206.
State 196 is now an ordinary working state and must not be interpreted as
halt when executing this table. Its SHA-256 is
\begin{quote}\small
\nolinkurl{7be1c776f1c5ac1db48b082b1880c710e0ce5722ea82fd66bc33e4ac561afe73}.
\end{quote}

The release verifier reconstructs the exact macro-word edit and checks its
postfix expression against the displayed gate $G'$. It reconstructs all
3,745 expanded primitive sites from the modified source grammar. All 1,127
decrement sites have both branches checked, giving $3,745+1,127=4,872$
uniform contracts. The largest address has length 38; the longest symbolic
trace uses 524 verifier steps. The unchanged 45 engine rows retain their
existing unbounded summaries. This particular grammar uses two extra return
scanner states.

An independent literal blank-tape execution reaches the exact initial
source frame after 59,164 steps, at head position $-2$ and state 75.
All 412 rows commute with the documented relabeling, and all 206 working
states are graph reachable. The raw and dense tables have identical startup
configurations under that relabeling.

A further reassociation and compatible-fold search gives a 203-state table
in \path{comparison/audit/integrated_gate203_dense.tm}, with start state 0,
working states $0,\ldots,202$, and halt label 203. Its SHA-256 is
\begin{quote}\small
\nolinkurl{c2da8e535f77ba8441ca443ad6f7e93262096eb8dc99373152a64d2774fb5e21}.
\end{quote}
The release verifier reconstructs its grammar, proves equality of the
expanded 3,745-instruction ordinal control graph with the initial integrated
source, and checks all 4,872 contracts. The largest source path is 39; the
longest symbolic trace uses 515 verifier steps. This table requires no extra
return-scanner states. All 406 rows pass the complete relabeling check and
all 203 working states are graph reachable. Its exact blank startup is
58,154 steps, with the same whole initial source frame, head $-2$, and root
75 as PA196.

Thus the source reduction has certified machine realizations. The best one
currently costs seven more physical states than the published PA196 table,
illustrating why source length and decoder-state count must be measured
separately.

\section{Why the apparent smaller candidates were rejected}
\label{sec:negative}
\subsection{A four-step loop caused by incompatible helper reuse}
The retained diagnostic
\path{verification/candidates/warm_physical_61_8047.tm}
has 195 working labels. It was produced by an optimistic combination of
reader folds and a longer return scan. It is not an accepted construction.

The release verifier starts it at the valid clean frame for source site
1111, with nine zero-valued live registers and 26 one-valued spare registers.
After 26 literal transitions the entire finite tape, head, and control enter
a cycle. The four control states are
\[
              1024\longrightarrow1056\longrightarrow1038
                   \longrightarrow1018\longrightarrow1024.
\]
The complete configuration repeats at step 30. Thus the claimed finite
source return does not occur. This is an exact finite-tape counterexample to
the proposed simulation, independently replayed without symbolic register
summaries. It is not merely a timeout at an arbitrarily chosen fuel limit.

The counterexample refutes the candidate's claimed compiler certificate.
It is not presented as a general decision procedure for consistency-machine
correctness, nor as a proof that every 195-state table fails the task.

\subsection{The return scan and transient zeros}
Another obstruction occurs even before incompatible folds are applied.
The unchanged return routine distinguishes the left tape margin using a
long run of zeros. Rebracketing changes the runs of zeros in source
addresses. Decrement preparation can introduce an additional transient zero,
so a bound on the unmodified address alone is insufficient.

For example, two optimistic grammar candidates with seven-zero address runs
pass most of their source contracts after adding one scan state, but fail
14 and 18 decrement contracts respectively. Adding a second scan state makes
their pure-reader versions pass all 4,959 contracts. Their compatible reader
folds then leave 197 working states. The apparent saving vanishes when both
the scan and simultaneous helper obligations are met.

Likewise, deleting the final state of the original return detector does not
establish a 195-state machine: the initial comparison address already
supplies a failing case. Source rebracketing, contextual macro entry, and
return-scan behavior have to be solved together.

\subsection{The precise scope of the searches}
The archive retains deterministic search scripts and selected receipts for
several families: direct identifications, mixed transition-row choices,
local detours, existing-helper substitutions, associativity changes,
exposure of closed macro bodies, altered acceptance conventions, physical
register allocations, axiom-head permutations, and equivalent logical
orientations. A separate integer program explores reusable macro phrases.
A solver time limit or an incumbent cost is not an optimality certificate.

Some particularly concrete negative results are:
\begin{itemize}
\item The original direct nonengine fold pass tested 33,150 candidates; the
mixed-row pass tested 28,440 candidates. Their protected engine and frame
conditions are stated in the scripts.
\item All 36 single transpositions of the nine physical register positions
were verified at the source boundary, and the corresponding direct-fold
passes tested 1,193,400 candidates without an accepted deletion.
\item All $8!=40,320$ register allocations fixing the original candidate
position were scored against compatible predecessor folds on the pure
202-state grammar. Each retained the same best six-state saving.
\item Several searches over hundreds of thousands of grammar mutations
returned valid or promising 196-state layouts, but no certified smaller
literal table. The results depend on their specified grammar and helper
families.
\end{itemize}
These are scoped computational results. They do not prove that the published
machine is locally optimal under every transformation, that a particular
source language is minimal, or that a sub-196 consistency machine cannot
exist.

\subsection{An append-zero initialization failure}
One experiment changes only the spare-register append action
\[
 (125,0):\quad 1L39\ \longmapsto\ 0L39.
\]
Even if later register contracts are reinterpreted to permit zero-valued
spares, this edit fails the required blank initialization. After exactly
76 steps the machine is in state 35 at head position 4, and the only
nonzero tape positions are $-3$ and 1. Its row on zero is $1R35$.
Every cell at or to the right of the head is therefore zero, and each next
step writes behind the new head and advances right in the same state.
This proves that the modified machine never reaches the intended initial
source frame. The release verifier reproduces the complete configuration
and the escaping-loop premise. The rejected edit is not included among
the positive machine constructions.

\section{Reproduction and proof boundaries}
\label{sec:reproduction}
\subsection{One command for the positive claims}
Extract the ZIP and, from its top-level directory, run
\begin{lstlisting}
python3 verify_release.py
python3 -O verify_release.py
\end{lstlisting}
The standard library is sufficient for these release checks. The script
replays the published and both new 196-state tables, verifies exact source
identity or the documented $h$-only source edit, checks all unbounded-engine
rows, verifies the six simultaneous folds and complete relabeling, replays
blank startup, reproduces the four-step rejected-candidate witness, and
repeats the 208-state graph census and neutral-helper checks.
It also reconstructs the shortened integrated-gate source, checks its
symbolic formula and all 4,872 source contracts for each realization,
verifies the 206- and 203-state relabelings and blank startups, and reproduces
the append-zero failure.
It writes its detailed outputs under \texttt{receipts/}.

The arithmetic behind the unbounded register summaries is in the pinned
upstream source snapshot under \texttt{sources/BusyBeaver/}. The independent
Python transcription is \texttt{verification/symbolic.py}. The full
comparison inputs and their byte hashes are in
\texttt{comparison/audit/inputs/}. The grammar optimization scripts are
research tools; their output must pass the release-style checks before it
is treated as a machine construction.

\subsection{Rebuilding this PDF}
On an ordinary TeX Live installation,
\begin{lstlisting}
pdflatex report.tex
pdflatex report.tex
\end{lstlisting}
rebuilds the PDF. The document uses Latin Modern, standard mathematical
packages, \texttt{booktabs}, \texttt{longtable}, \texttt{xurl}, and
\texttt{hyperref}. The included build script also documents the local-format
fallback used in the execution environment. All TeX inputs are in the ZIP.

\subsection{What has and has not been proved here}
The exact finite table comparisons, graph identifications, literal startups,
source-site replay calculations, and failed-cycle witness are executable
results. The mathematical interpretation and integrated-gate proof provide
explicit uniform constructions and explain their internalization in PA.
They are separate from a completed kernel-checked formalization. The
upstream Lean declarations retain their original claims and identities;
no new formal theorem is attributed to them.

The final state-count conclusion is limited and explicit: the supplied
archive contains verified new 196-state variants and a smaller logical
source program with a verified 203-state implementation. It does not supply
the requested verified machine with fewer than 196 working states.

\appendix
\input{reduction196/grammar_milp_note.tex}
\section{Complete transition table of the rebracketed machine}
\label{app:table}
The start state is 0, the blank symbol is 0, and $H$ denotes halt label 196.
An entry $wDq$ writes $w$, moves one cell in direction $D$, and enters $q$.
The three blocks below cover every working state and both scanned symbols.
\input{transition_table.tex}
\end{document}
